\section{Casos criticos de prueba}

\subsection{Descripcion}

Este documento define el inventario minimo de escenarios de negocio que no deben romperse. Su objetivo es proteger la justicia competitiva, la consistencia operativa y la confianza del sistema sobre los puntos de mayor riesgo del producto.

\subsection{Alcance}

Aplica a activacion, picks, lock, resultados, scoring, leaderboard, administracion y contingencias relevantes para v1.

\subsection{Definiciones}

\begin{itemize}
  \item \textbf{Caso critico}: escenario cuya falla afecta justicia, operacion o confianza del sistema.
  \item \textbf{Resultado esperado}: comportamiento verificable que debe cumplirse sin ambiguedad.
\end{itemize}

\subsection{Reglas}

\begin{itemize}
  \item Los casos criticos se definen por impacto funcional, no solo por complejidad tecnica.
  \item Toda regla critica del juego debe tener al menos un caso de prueba asociado.
  \item Los casos deben poder mapearse claramente al documento responsable de la regla que protegen.
\end{itemize}

\subsubsection{Inventario minimo de casos criticos}

\begin{longtable}{p{0.06\textwidth}p{0.34\textwidth}p{0.22\textwidth}p{0.24\textwidth}}
\toprule
\textbf{ID} & \textbf{Escenario} & \textbf{Area principal} & \textbf{Resultado esperado} \\
\midrule
CT-01 & usuario \texttt{PENDIENTE} intenta guardar pick & activacion y picks & operacion rechazada por estado funcional \\
CT-02 & usuario \texttt{ACTIVO} guarda pick valido antes del lock & picks & pick persistido correctamente \\
CT-03 & usuario intenta editar pick despues del lock & locks & operacion rechazada y pick permanece bloqueado \\
CT-04 & partido llega al lock sin prediccion guardada & locks y scoring & se interpreta \texttt{SIN\_PICK} y no se fabrica valor \\
CT-05 & prediccion especial intenta editarse despues del cierre del primer partido & predicciones especiales & operacion rechazada por lock especial \\
CT-06 & resultado oficial llega y se calcula scoring base & scoring & puntaje por outcome y exacto coincide con la regla \\
CT-07 & partido de fase con multiplicador aplica precision exacta & scoring & total calculado sin redondeo interno arbitrario \\
CT-08 & partido cancelado sin resultado utilizable & Edge Cases & no genera puntaje para ningun usuario \\
CT-09 & resultado corregido dispara recalculo & resultados y leaderboard & ranking y puntajes se actualizan desde verdad vigente \\
CT-10 & override manual vigente recibe nueva sincronizacion automatica & integraciones & el valor manual no es sobrescrito automaticamente \\
CT-11 & admin aprueba comprobante valido & activacion admin & usuario pasa a \texttt{ACTIVO} con auditoria \\
CT-12 & admin rechaza comprobante & activacion admin & usuario pasa a \texttt{RECHAZADO} con trazabilidad \\
CT-13 & accion administrativa sin permisos suficientes & seguridad & operacion rechazada y sin cambio de estado \\
CT-14 & archivado del torneo bloquea operacion competitiva ordinaria & admin y operacion & competencia queda en modo solo lectura \\
\bottomrule
\end{longtable}

\subsubsection{Casos criticos adicionales recomendados}

\begin{itemize}
  \item cambio de kickoff antes del lock recalcula correctamente la ventana abierta;
  \item postergacion despues del lock preserva picks ya bloqueados salvo accion extraordinaria auditada;
  \item leaderboard muestra total acumulado coherente con desglose cuando ese desglose se exponga;
  \item mensaje funcional por resultado corregido aparece en la superficie definida sin exponer detalle operativo excesivo.
\end{itemize}

\subsubsection{Asignacion por nivel de prueba}

\begin{itemize}
  \item \texttt{CT-01} a \texttt{CT-07}: backend e integracion como minimo, con apoyo end to end en flujos principales.
  \item \texttt{CT-08} a \texttt{CT-10}: integracion y pruebas operativas por su dependencia de estados no ordinarios.
  \item \texttt{CT-11} a \texttt{CT-14}: integracion y end to end administrativo.
\end{itemize}

\subsection{Edge Cases}

\begin{itemize}
  \item Un mismo caso critico puede requerir mas de una prueba si existe riesgo tanto en backend como en frontend.
  \item Algunos casos criticos dependen de datos temporales o scheduler; en esos casos la automatizacion puede requerir dobles o fixtures controlados.
  \item Un caso visible en UI no se considera cubierto si falta verificacion del comportamiento real del backend.
\end{itemize}

\subsection{Invariantes}

\begin{itemize}
  \item Ninguna regla critica del juego queda sin un caso de prueba asociado.
  \item Todo caso critico debe tener resultado esperado verificable.
  \item Los escenarios de lock, scoring y correccion de resultados mantienen prioridad alta en cualquier iteracion.
\end{itemize}

\subsection{Preguntas abiertas}

\begin{itemize}
  \item Confirmar si el proyecto quiere formalizar una matriz trazable completa entre cada caso critico y su tipo de prueba exacto por repositorio.
\end{itemize}

\subsection{Decisiones pendientes}

\begin{itemize}
  \item \texttt{Decision pendiente} \texttt{No bloqueante}: decidir si se agregara una codificacion oficial de casos de prueba para rastreo transversal entre repositorios.
  \item \texttt{Decision pendiente} \texttt{No bloqueante}: definir si v1 exigira cobertura automatizada obligatoria para todos los casos de este inventario o si algunos podran validarse manualmente en salida controlada.
\end{itemize}
